% CN 1.16 Exercises (p.28--35)

% ======================================================================
\begin{exercise}{1}
    Why does the set difference only satisfy (1.9) $\bigl(|B\setminus A| = |B|-|A|\bigr)$
    when $A\subseteq B$? How would you modify (1.9) so that it covers any set
    difference $B\setminus A$? What extra information about the sets would you
    need, in order to determine $|B\setminus A|$?
\end{exercise}

% ======================================================================
\begin{exercise}{2}
    In C, \texttt{int monthNumber;} declares a variable of type \texttt{int},
    and \texttt{char monthName[10];} declares a string of at most 9 characters.
    Let \texttt{Int} be the set of possible values for a variable of type
    \texttt{int}, and \texttt{String} the set of possible values for a string of
    at most 9 letters.
    \begin{enumerate}[label=(\alph*)]
        \item The statement
\begin{verbatim}
struct aNewType {
    int year;
    char monthName[10];
} monthBothWays;
\end{verbatim}
        creates a type for objects consisting of any int followed by any string.
        Using \texttt{Int}, \texttt{String} and a standard set operation, what
        set is represented by \texttt{aNewType}?
        \item The statement
\begin{verbatim}
union anotherNewType {
    int year;
    char monthName[10];
} monthEitherWay;
\end{verbatim}
        creates a type for objects that can be either an int or a string (only
        one at a time). What set is represented by \texttt{anotherNewType}?
    \end{enumerate}
\end{exercise}

% ======================================================================
\begin{exercise}{3}
    Suppose $A$ and $B$ are subsets of some universal set $U$.
    \begin{enumerate}[label=(\alph*)]
        \item If $A\subseteq B$, what is $A\cup B$?
        \item If $A\not\subseteq B$, what can you say about $A\cup B$?
        \item Complete the following: $A\subseteq B$ if and only if
              $A\cup B = \underline{\hspace{2cm}}$.
        \item Devise another equivalent condition for $A\subseteq B$ involving
              intersection instead of union.
        \item Devise another equivalent condition for $A\subseteq B$ involving
              set difference.
    \end{enumerate}
\end{exercise}

% ======================================================================
\begin{exercise}{4}
    Consider diagrams of the following type (e.g.\ for $\{a,b\}$: $\emptyset$ at
    the bottom, $\{a\},\{b\}$ in the middle, $\{a,b\}$ at the top). The
    requirements are:
    \begin{itemize}
        \item Sets are written in their usual textual form (not Venn diagrams).
        \item Smaller sets are lower on the page than larger sets; sets of the
              same size are on the same horizontal level.
        \item If $X\subset Y$ and $|Y|=|X|+1$ then there is an arrow from $X$ to $Y$.
        \item The diagram is as neat and symmetric as possible.
    \end{itemize}
    \begin{enumerate}[label=(\alph*)]
        \item Draw a diagram of this type that shows all subsets of $\{a,b,c\}$
              and the subset relation among them.
        \item For every pair $X,Y$ with $X\subset Y$ and $|Y|=|X|+1$, label the
              corresponding arrow by the sole member of $Y\setminus X$.
        \item If you could draw in three dimensions, how could you draw this
              diagram in a natural way?
        \item For a set of $n$ elements, how many sets and how many arrows does a
              diagram of this type have?
        \item For each element of the $n$-element set in (d), how many arrows are
              labelled by that element?
        \item Suppose $X\subseteq Y$. How many directed paths (along arrows,
              forwards only) are there from $X$ to $Y$? Give an expression.
        \item With $X,Y$ as in (f), what is the maximum number of mutually
              internally disjoint paths from $X$ to $Y$? (Two paths are
              internally disjoint if no internal set on either path appears
              anywhere on the other path.)
        \item Explain how to use the diagram to find, for any two sets in it,
              their union and intersection.
    \end{enumerate}
\end{exercise}

% ======================================================================
\begin{exercise}{5}
    Given $U=\{e_1,e_2,\dots,e_n\}$, the \emph{characteristic string} of a subset
    $A\subseteq U$ is the $n$-bit string $b_1b_2\cdots b_n$ where
    \[
        b_i = \begin{cases} 1, & e_i\in A;\\ 0, & e_i\notin A.\end{cases}
    \]
    Write down the characteristic string of each of the subsets of a set of three
    elements. List them, one above the other, so that each differs from the one
    above it in just one bit. See if you can extend this to subsets of a set of
    $n$ elements.
\end{exercise}

% ======================================================================
\begin{exercise}{6}
    Two intersecting sets $A,B$ divide the plane into four basic regions:
    $\overline{A\cup B}$, $A\cap B$, $A\setminus B$, $B\setminus A$. Three sets
    can divide the plane into eight regions.
    \begin{enumerate}[label=(\alph*)]
        \item Label each basic region of a Venn diagram for three sets $A,B,C$
              with appropriate intersections of sets.
    \end{enumerate}
    A \emph{general Venn diagram} for $n$ sets $A_1,\dots,A_n$ is one in which
    every intersection of the form
    \[
        (A_1 \text{ or } \overline{A_1})\cap(A_2\text{ or }\overline{A_2})\cap\cdots
        \cap(A_n\text{ or }\overline{A_n})
    \]
    corresponds to one of the basic regions, so there are $2^n$ basic regions.
    \begin{enumerate}[label=(\alph*),start=2]
        \item What is the maximum number of sets for which a general Venn diagram
              can be drawn in which all the sets are circles of the same size?
        \item Draw a general Venn diagram in the plane for four sets. You can use
              closed curves other than circles.
        \item How could you draw a three-dimensional general Venn diagram using
              four sets, each represented as a sphere?
    \end{enumerate}
\end{exercise}

% ======================================================================
\begin{exercise}{7}
    Prove that if $A\subseteq B$ then $\power{A}\subseteq\power{B}$.
\end{exercise}

% ======================================================================
\begin{exercise}{8}
    Consider the sequence of binomial coefficients
    $\binom n0,\binom n1,\binom n2,\dots,\binom n{n-1},\binom nn$.
    \begin{enumerate}[label=(\alph*)]
        \item Using one of the formulas for $\binom nr$, prove that these
              binomial coefficients increase as $r$ goes from $0$ to
              $\lfloor n/2\rfloor$ and then decrease as $r$ goes from
              $\lceil n/2\rceil$ to $n$. (In the odd case, what happens when $r$
              goes from $(n-1)/2$ to $(n+1)/2$?)
        \item Prove that, for every integer $r$ with $1\le r\le n-1$,
              \[
                  \binom nr^2 > \binom n{r-1}\binom n{r+1}.
              \]
              (The sequence is \emph{strictly log-concave}.)
        \item You choose an ordered pair of subsets of size $r$ of a set of size
              $n$ (no other restriction). Your friend chooses one subset of size
              $r-1$ and another of size $r+1$. Who has more options? Does the
              answer depend on $n$ and $r$? Is there any situation where the
              numbers of choices are the same?
    \end{enumerate}
\end{exercise}

% ======================================================================
\begin{exercise}{9}
    Express $|A\cup B|$ in terms of $|A|,|B|,|A\cap B|$.
\end{exercise}

\begin{exercise}{10}
    Express $|A\cap B|$ in terms of $|A|,|B|,|A\cup B|$.
\end{exercise}

\begin{exercise}{11}
    Express $|A\cup B\cup C|$ in terms of
    $|A|,|B|,|C|,|A\cap B|,|A\cap C|,|B\cap C|,|A\cap B\cap C|$.
\end{exercise}

\begin{exercise}{12}
    Express $|A\cap B\cap C|$ in terms of
    $|A|,|B|,|C|,|A\cup B|,|A\cup C|,|B\cup C|,|A\cup B\cup C|$.
\end{exercise}

% ======================================================================
\begin{exercise}{13}
    Suppose we have sets $A_1,\dots,A_n$. For each $k$, let
    \[
        i_k := \text{sum of all sizes of intersections of $k$ of the sets.}
    \]
    E.g.\ for $n=3$: $i_1=|A_1|+|A_2|+|A_3|$,
    $i_2=|A_1\cap A_2|+|A_1\cap A_3|+|A_2\cap A_3|$, $i_3=|A_1\cap A_2\cap A_3|$.

    Rewrite your expressions for $|A\cup B|$ and $|A\cup B\cup C|$ in terms of
    $i_1,i_2,\dots$. Then write out a general expression for
    $|A_1\cup\cdots\cup A_n|$. Pay careful attention to the sign of the last term.
\end{exercise}

% ======================================================================
\begin{exercise}{14}
    Now let
    \[
        u_k := \text{sum of all sizes of unions of $k$ of the sets.}
    \]
    Rewrite your expressions for $|A\cap B|$ and $|A\cap B\cap C|$ in terms of
    $u_1,u_2,\dots$. Then write out a general expression for
    $|A_1\cap\cdots\cap A_n|$. Pay careful attention to the sign of the last term.
\end{exercise}

% ======================================================================
\begin{exercise}{15}
    Express $|A\triangle B|$ in terms of the sizes of some other sets in three
    different ways.
\end{exercise}

\begin{exercise}{16}
    Draw a Venn diagram for general sets $A,B,C$ and shade the region(s) that
    form $A\triangle B\triangle C$.
\end{exercise}

\begin{exercise}{17}
    Prove that, for any sets $A$ and $B$,
    \[
        (A\cup B)\triangle(A\cap B) = A\triangle B.
    \]
\end{exercise}

\begin{exercise}{18}
    Suppose $A_1,\dots,A_n$ are sets. How would you characterise
    $A_1\triangle A_2\triangle\cdots\triangle A_n$? More specifically, its members
    are those that satisfy some condition on how many of the sets they belong to.
    What is that condition?
\end{exercise}

% ======================================================================
\begin{exercise}{19}
    If $A$ and $B$ are sets, does $\overline{A\times B}$ equal
    $\overline A\times\overline B$? If so, prove it. If not: give an example when
    it does not hold; characterise those cases when it does hold; and determine
    what can be said about the relationship between them.
\end{exercise}

\begin{exercise}{20}
    Write down all partitions of $\{1,2,3,4\}$ into three parts.
\end{exercise}

\begin{exercise}{21}
    Suppose $A,B,C\subseteq U$ are nonempty sets such that all basic regions in
    their general Venn diagram are nonempty (see Exercise 6). Then $A,B,C$ define
    a partition of $U$ into eight parts. What are those parts? Express each part
    in terms of one or more of $A,B,C$ using set operations.
\end{exercise}

\begin{exercise}{22}
    Does $\R$ have a partition in which all parts are
    \begin{enumerate}[label=(\alph*)]
        \item open intervals?
        \item half-open half-closed intervals?
        \item closed intervals?
    \end{enumerate}
\end{exercise}
